\documentclass[10pt,a4paper]{amsart}
\usepackage[margin=19mm]{geometry}
\usepackage{amsmath,amssymb,mathtools}
\usepackage[scr=rsfs,cal=euler]{mathalfa}
\usepackage{xfrac}
\usepackage[unicode,hidelinks,bookmarks=false]{hyperref}
\newcommand{\ie}{i.e.\ }
\newcommand{\cf}{cf.\ }
\newcommand{\resp}{respectively\ }
\newcommand{\Subsection}{\subsection}
\providecommand{\nobreakdash}{\nobreak\hbox{-}}
\setlength{\emergencystretch}{2em}
\multlinegap0pt
\usepackage{upref}
\usepackage{cleveref}
\usepackage{mathtools}
\usepackage{tikz-cd}
\usepackage{thmtools}
\newtheorem{theorem}{Theorem}[section]
\newtheorem{lemma}[theorem]{Lemma}
\newtheorem{corollary}[theorem]{Corollary}
\newtheorem{remark}[theorem]{Remark}
\newtheorem{example}[theorem]{Example}
\newtheorem{problem}[theorem]{Problem}
\newtheorem{definition}[theorem]{Definition}
\newtheorem{proposition}[theorem]{Proposition}
\newcommand{\Sym}{\operatorname{Sym}}
\newcommand{\Aut}{\operatorname{Aut}}
\newcommand{\Fix}{\operatorname{Fix}}
\newcommand{\tr}{\operatorname{tr}}
\newcommand{\PSL}{\operatorname{PSL}}
\newcommand{\PGL}{\operatorname{PGL}}
\newcommand{\PGammaL}{\operatorname{P\Gamma L}}
\newcommand{\PG}{\operatorname{PG}}
\begin{document}
\title[Characterization of Weak EKR Groups and Intersection Densiti]{Characterization of Weak EKR Groups and Intersection Densities with Prescribed Point Stabilizers}
\author{Boštjan Frelih, Ademir Hujdurović, Klavdija Kutnar}
\date{}
\maketitle
\noindent\emph{Source and licence.} Characterization of Weak EKR Groups and Intersection Densities with Prescribed Point Stabilizers. Version arXiv:2608.18594v1 (CC BY 4.0). This material is distributed under the Creative Commons Attribution 4.0 International licence, http://creativecommons.org/licenses/by/4.0/, as recorded in the arXiv OAI licence metadata for this exact version and shown on the abstract page https://arxiv.org/abs/2608.18594v1. Full rights notice: licenses/arxiv-2608.18594-CC-BY-4.0.txt.\par\smallskip
\noindent\textit{Editorial scope.} Counted unit: the original \texttt{theorem} environment \texttt{thm:normal-extension} at source lines 268--280 together with its complete proof at lines 282--356; the delivered document keeps source lines 185--357 so that every referenced label and every citation resolves inside it. Native editable source is the paper's own concatenated TeX; the AMS wrapper is editorial formatting only. No publisher class, upstream script or unrelated artwork is redistributed.
\par\smallskip\noindent\textit{Reference note.} This article ships no compiled bibliography (biblatex); the entries delivered here are composed from the author's own BibTeX (.bib) fields, with no field invented and no entry dropped.
\section{Characterization of weak EKR groups}



Throughout this section, all groups are finite. For subgroups and
conjugates we use the convention
\[
H^g=g^{-1}Hg,
\]
and we write
\[
\mathcal{U}_G(H):=\bigcup_{g\in G}H^g
\]
for the union of all conjugates of \(H\) in \(G\).

For \(H\leq G\), let \(\rho(G,H)\) denote the intersection density of
the action of \(G\) on the right cosets of \(H\). Thus
\[
\rho(G,H)
 =
 \frac{1}{|H|}
 \max\left\{
 |\mathcal{F}|:
 \mathcal{F}\subseteq G,\ 
 \mathcal{F}^{-1}\mathcal{F}\subseteq \mathcal{U}_G (H)
 \right\}.
\]
In particular, \(G\) has the weak EKR property if and only if
\[
\rho(G,H)=1
\qquad\text{for every }H\leq G.
\]

\begin{definition}[Local density]
\label{def:local-defect}
Let \(H\leq N\trianglelefteq G\). Define the {local density of \(H\) in \(G\) relative to \(N\)} by
\[
\rho_N(G,H)
 :=\frac{1}{|H|}
 \max\left\{
 |A|:
 A\subseteq N,\ 
 A^{-1}A\subseteq \mathcal{U}_G(H)
 \right\}.
\]
 

\end{definition}

Observe that $\rho_N(G,H)\geq 1$ since we can always take \(A=H\) in the definition of
\(\rho_N(G,H)\). It is also clear that $\rho_N(G,H)\leq \rho(G,H)$, as in $\rho(G,H)$ the maximum is taken over all subsets of $G$, while in $\rho_N(G,H)$ the maximum is taken only over subsets of $N$. In the next lemma we prove that we actually have the equality.

\begin{lemma}
\label{lem:local-defect-realization}
Let \(N\trianglelefteq G\) and let \(H\leq N\). Then
\(
\rho_N(G,H)=\rho(G,H).
\)
\end{lemma}

\begin{proof}
Let \(\mathcal{F}\subseteq G\) be a maximum intersecting  set in the action of \(G\)
on cosets of \(H\). By the definition of intersecting set it follows that
\(
x^{-1}y\in \mathcal{U}_G(H)\)
for all \(x,y\in\mathcal{F}.
\)
Since \(H\leq N\) and \(N\trianglelefteq G\), we have
\(
\mathcal{U}_G(H)\subseteq N.
\)
It follows that \(x^{-1}y\in N\) for all \(x,y\in\mathcal{F}\), and
hence all elements of \(\mathcal{F}\) lie in the same right coset of
\(N\).

Without loss of generality we may assume that $1\in \mathcal{F}$, therefore we may assume that
\(\mathcal{F}\subseteq N\). By the definition of $\rho_N(G,H)$ it follows that $\rho_N(G,H)\geq \frac{|\mathcal{F}|}{|H|}$. Since $\mathcal{F}$ is a maximum intersecting set, we have 
$\rho(G,H)=\frac{|\mathcal{F}|}{|H|}$. 
This shows that $\rho_N(G,H)\geq \rho(G,H)$ hence we obtain the equality.
\end{proof}

We will now present the characterization of weak EKR groups.


\begin{theorem}
\label{thm:normal-extension}
Let \(N\trianglelefteq G\). Then \(G\) has the weak EKR property if
and only if the following two conditions hold:
\begin{enumerate}
\item \(G/N\) has the weak EKR property;
\item
\(
\rho_N(G,H)=1
\quad\text{for every }H\leq N.
\)
\end{enumerate}
\end{theorem}

\begin{proof}
Suppose first that \(G\) has the weak EKR property. The weak EKR
property is inherited by quotient groups by \cite[Lemma 4]{BardestaniMallahiKarai2015}, and hence \(G/N\) also has the
weak EKR property. Since $G$ has the weak EKR property, then $\rho(G,H)=1$ for every subgroup $H\leq G$, in particular also for every $H\leq N$. By Lemma~\ref{lem:local-defect-realization} it follows that $\rho_N(G,H)=1$, for every $H\leq N$.


Conversely, suppose that \(G/N\) has the weak EKR property and that
\(
\rho_N(G,K)=1\)
for every \(K\leq N.
\)
Let \(H\leq G\), and let \(\mathcal{F}\subseteq G\) be intersecting set in
the action of \(G\) on the cosets of \(H\). Let
\(
K:=H\cap N
\)
and let
\(
\pi:G\rightarrow G/N
\)
be the quotient map.

It is straightforward to check that the set \(\pi(\mathcal{F})\) is intersecting in the action
of \(
G/N \) on  cosets of $HN/N$.
Since by the assumption \(G/N\) has the weak EKR property, it follows that $\rho(G/N,HN/N)=1$, hence
\begin{equation}    
|\pi(\mathcal{F})|
 \leq
 |HN/N|
 =
 \frac{|H|}{|H\cap N|}
 =
 \frac{|H|}{|K|}.
\label{eq:image-bound}
\end{equation}

Let $f\in \mathcal{F}$ and consider the intersection of $\mathcal{F}$ with the coset $fN$.  Let $S\subseteq N$ such that $\mathcal{F} \cap fN=fS$. If \(a,b\in S\), then \(fa,fb\in
\mathcal{F}\), so
\[
a^{-1}b=(fa)^{-1}(fb)\in H^g
\]
for some \(g\in G\). Since \(a^{-1}b\in N\), the normality of \(N\) gives
\[
a^{-1}b
 \in
 N\cap H^g
 =
 (N\cap H)^g
 =
 K^g.
\]
Consequently,
\(
S^{-1}S\subseteq \mathcal{U}_G(K).
\)
Since \(\rho_N(G,K)=1\), it follows that
\(
|S|
 \leq |K|.
\) 
This shows that $|\mathcal{F}\cap fN|\leq |K|$.
Together with \eqref{eq:image-bound}, it follows that
\[
|\mathcal{F}|
 \leq
 |K|\,|\pi(\mathcal{F})|
 \leq
 |K|\frac{|H|}{|K|}
 =
 |H|.
\]
Therefore the action of \(G\) on the cosets of \(H\) has the EKR property. Since
\(H\leq G\) was arbitrary, \(G\) has the weak EKR property.
\end{proof}


\begin{thebibliography}{99}
\bibitem[]{BardestaniMallahiKarai2015} Bardestani, Mohammad; Mallahi-Karai, Keivan. On the {Erd\H{o}s--Ko--Rado} property for finite groups. Journal of Algebraic Combinatorics 42 (2015) 111--128
\end{thebibliography}
\end{document}
